\section{Experimental Setup}
\label{sec:experiment}

\subsection{Hardware}
\label{sec:hardware}

The actuator under test is a six-pole (hexapole) electromagnetic actuator
without a ferromagnetic yoke. Each pole consists of a soft-iron core wound
with $\Nc = 50$ turns of copper wire. The pole tips are machined to a
curvature radius of $\rtip = \SI{5}{\um}$.

The excitation signal is generated by a 16-bit DAC (zero offset 32768, range
$\pm\SI{10}{\volt}$, resolution 65536) and amplified by a transconductance
amplifier with gain $\kA = \SI{0.3614}{\ampere/\volt}$ to drive the coil.
The magnetic field is measured by an EQ-730L linear Hall effect
sensor~\cite{EQ730L} with a sensitivity of $\SH = \SI{130}{\volt/\tesla}$
and a linear range of $\pm\SI{15}{\milli\tesla}$.

Data acquisition is performed at a sampling rate of
$f_s = \SI{20}{\kilo\hertz}$ using V8-format binary files, with both the
DAC output voltage and the Hall sensor voltage recorded simultaneously.

\subsection{Measurement Protocol}
\label{sec:protocol}

The experiment is organized as a two-dimensional grid of
22~distances $\times$ 15~frequencies, yielding 330 individual measurements.
At each distance, a single-frequency sinusoidal excitation is applied to the
coil, and the steady-state Hall sensor response is recorded.

The 22 distances span from \SI{10}{\um} to \SI{2000}{\um}:
\begin{equation*}
x \in \{10,\; 30,\; 50,\; 100,\; 150,\; 200,\; 250,\; 300,\; 350,\; 400,\;
450,\; 500,\; 600,\; 700,\; 800,\; 900,\; 1000,\; 1200,\; 1400,\; 1600,\;
1800,\; 2000\} \;\si{\um}.
\end{equation*}

The 15 excitation frequencies are:
\begin{equation*}
f \in \{0.1,\; 1,\; 10,\; 50,\; 100,\; 200,\; 500,\; 1000,\; 1200,\; 1250,\;
1500,\; 1600,\; 2000,\; 2500,\; 3000\} \;\si{\hertz}.
\end{equation*}

\subsection{Data Processing}
\label{sec:data-processing}

Each recorded time series undergoes the following processing pipeline:

\begin{enumerate}
\item \textbf{Steady-state detection}: A conservative relative-threshold
algorithm (threshold 0.2\%, 3 consecutive periods) identifies the onset of
steady-state oscillation. Low-frequency files that do not reach steady state
use the last available periods as a fallback.

\item \textbf{Super-period truncation}: The steady-state segment is truncated
to an integer number of super-periods---the minimum number of samples
containing an exact integer number of excitation cycles---to eliminate spectral
leakage in the subsequent FFT. For non-integer-period frequencies, a fallback
to the nearest integer period is used.

\item \textbf{FFT}: The transfer function magnitude is computed as
$|H| = |V_{\mathrm{m,FFT}}(f)| / |V_{\mathrm{DA,FFT}}(f)|$
at the excitation frequency bin.
\end{enumerate}

\subsection{Data Quality}
\label{sec:data-quality}

All 330 out of 330 distance--frequency combinations yielded valid transfer
function magnitudes. \Cref{fig:lissajous} shows representative Lissajous
figures at $f = \SI{0.1}{\hertz}$ (the lowest frequency, where nonlinearity
or drift would be most apparent), confirming that the system operates in the
linear regime across all distances.

\begin{figure}[htbp]
\centering
\includegraphics[width=0.48\textwidth]{../results/Charge_NTU/figures/Charge_Lissajous_0p1Hz.png}%
\hfill
\includegraphics[width=0.48\textwidth]{../results/Charge_NTU/figures/Charge_Lissajous_0p1Hz_Normalized.png}
\caption{Lissajous figures at $f = \SI{0.1}{\hertz}$ for all 22 distances.
Left: raw $\Vm$ vs.\ $\Vda$. Right: normalized.
The elliptical traces confirm linear system behavior.}
\label{fig:lissajous}
\end{figure}
